\begin{textsample}{POS Dim 1 – vem\_ed\_09 – Score 7.00 – vem\_ed\_09/t037.txt}  \label{ex:f1_pos_087}
Professores avaliam \textbf{intercâmbios} virtuais Na pesquisa de percepção finalizada em 30 de julho de 2021, 30 professores orientadores de PCIs responderam ao questionário online.
Destes, apenas 8 estavam participando pela primeira vez de um PCI; 22 já \textbf{são} veteranos em Intercâmbios Virtuais.
Isso ocorre porque, ao concluir o projeto, os docentes se sentem muito motivados para continuar com mais um PCI no semestre seguinte.
Dos 28 projetos realizados, 20 \textbf{foram} em inglês, 6 em espanhol e 2 em português (com o Instituto Politécnico de Viseu, em colaboração \textbf{iniciada} no primeiro semestre de 2021).
Desde 2013, quando o primeiro Projeto Colaborativo Internacional (PCI) \textbf{foi} realizado, o \textbf{número} de participantes vem aumentando consideravelmente, como \textbf{é} \textbf{possível} verificar no gráfico “Evolução de participantes em PCIs”.
Houve um ligeiro declínio no primeiro semestre de 2020, devido à interrupção das atividades no \textbf{início} da pandemia.
Somando \textbf{todos} os PCIs nesses últimos oito anos nas Fatecs, aproximadamente 3.900 \textbf{pessoas} realizaram Intercâmbios Virtuais.
Pesquisa com professores dos PCIs (julho/2021) 25 professores consideram \textbf{que} a \textbf{interação} com seus alunos melhorou com os PCIs, 5 avaliam \textbf{que} ficou igual 29 professores avaliam a \textbf{interação} com o colega estrangeiro “ótima” ou “boa”, somente 1 considera regular

% matched lemmas: iniciar, interação, intercâmbio, início, número, pessoa, possível, que, ser, todo
\end{textsample}
